\section{Kernel and hull restrictions}\label{sec:restrictions}

We derive necessary conditions from the kernel and image of the extension
map. Throughout this section, \(S\) is a stabilizer with
\(S\cap V_T=R\), \(S^\perp\cap W=C\), and extension map \(M\).
Every nonzero kernel element represents a logical operator supported on
\(T\).
There are two kernels to distinguish. The kernel \(K=\ker M\) records
logical operators supported on \(T\). The kernel \(N=\ker\tau\)
records stabilizer words supported on the complementary coordinates:
their projected local class is zero, and an element of \(R\) removes
the remaining local representative. The hull identities below relate
these two kinds of local information.

\begin{proposition}[Kernel and hull identities]\label{prop:kernel}
For an extension, put \(K=\ker M\), \(q=\dim K\), and let \(t\) be the
rank of the symplectic form restricted to \(K\). Then
\begin{align}
 \rank\Gamma&=a-2q+t,&
 h&=b-a+2q-t,\label{eq:hull}\\
 N:=\ker\tau&\subseteq\rad(P),&
 \dim N&=b-a+q.\label{eq:N}
\end{align}
Every nonzero \(e\in K\) has \(\cost(e)\geq d(S)\), even without
purity. Moreover \(q-t/2\leq k\).
\end{proposition}
\begin{proof}
The adjoint map has image \(K^\perp_E\), of dimension \(a-q\).
To see this directly, \(e\) annihilates \(\im\tau\) if and only
if \(\langle e,\tau(x)\rangle=0\) for every \(x\in P\), or
equivalently \(Me=0\). Nondegeneracy of \(E\) then gives
\(\im\tau=K^\perp_E\).
The radical of its restricted form is \(K\cap K^\perp\), of dimension
\(q-t\) by rank--nullity on the restricted form. Pullback along the surjection \(P\to K^\perp\) leaves
the form rank unchanged, giving \(a-q-(q-t)\). Subtracting from \(b\)
gives \eqref{eq:hull}. The map \(\tau\) has rank \(a-q\), giving
\eqref{eq:N}; a vector in its kernel pairs trivially with all of \(P\).

For \(e\ne0\) in \(K\), a prefix representative with zero suffix lies
in \(S^\perp\). It cannot lie in \(S\), because \(S\cap V_T=R\).
Its minimum weight is at least the logical distance. The space
\(K\) embeds in the \(2k\)-dimensional nondegenerate logical space
\(S^\perp/S\). A subspace of dimension \(q\) and form rank \(t\)
has an isotropic subspace of dimension \(q-t/2\); that dimension is at
most \(k\).
\end{proof}

The local centralizer \(H=S^\perp\cap V_T\) satisfies \(H\cap S=R\)
and \(H/R=K\): a local class lies in \(\ker M\) exactly when it has a
representative with identity suffix.

\begin{proposition}[Kernel projection bound]\label{prop:griesmer}
Suppose \(R=\langle Z^m\rangle\) and \(d>\lfloor m/2\rfloor\).
The physical \(X\)-projection maps \(K\) injectively to an even binary
linear code of length \(m\), dimension \(q\), and distance at least
\(\Delta=2(d-\lfloor m/2\rfloor)\). Consequently,
\begin{equation}\label{eq:griesmer}
 m\geq\sum_{i=0}^{q-1}\left\lceil\frac{\Delta}{2^i}\right\rceil.
\end{equation}
\end{proposition}
\begin{proof}
Adding \(Z^m\) leaves the \(X\)-mask unchanged, so this projection is
well defined on \(E\). Commutation with \(Z^m\) makes its weight
\(w\) even. If \(z\) counts \(Z\)-letters off that mask, then
\[
 \cost(e)=w+\min(z,m-w-z)\leq\lfloor m/2\rfloor+w/2.
\]
A nonzero class represented by a \(Z\)-only word has minimum weight at most \(\lfloor m/2\rfloor\),
so cannot lie in \(K\). This proves injectivity, and
\(\cost(e)\geq d\) gives \(w\geq\Delta\).
The last inequality is the binary Griesmer bound applied to that projected
code. For completeness, choose a minimum-weight binary word of weight
\(w\). Puncturing its support leaves a code of dimension one less and
distance at least \(\lceil w/2\rceil\): the two representatives of a
nontrivial quotient class have weight sum equal to \(w\) plus twice
their common suffix weight, and each has weight at least \(w\).
Only zero and the chosen word can lie wholly in that support. Induction
gives the stated sum.
\end{proof}

For \(m=8\) and \(d=6\), the bound gives \(q\leq4\).
Lemma~\ref{lem:qthree} sharpens this to three using parity and the
logical symplectic dimension.
